\section{Scalar tests and incidence averaging}\label{sec:tests}

Keep the finite state set $S$, the integer $n\ge2$, and the graph
$G=G_m$ from the continuation construction in \Cref{sec:walks}.
For an incidence $e:x\to y$, we will prescribe the bump $b_x$ on
states outside $E_e$ and the negative bump $-b_y$ on states outside
$E_{\bar e}$. The bump bounds handle pairs within one set and cross pairs
at distance at least $2L$. The exceptional-direction lemma handles closer
cross pairs, making the combined prescribed values $1$-Lipschitz.

\begin{lemma}\label{lem:edge-test}
For every incidence $e:x\to y$ of $G$, there is a $1$-Lipschitz
function $f_e:M_p\to[-L,L]$ such that
\[
 f_e(j)=b_x(j)\quad(j\in S\setminus E_e),\qquad
 f_e(j')=-b_y(j')\quad(j'\in S\setminus E_{\bar e}).
\]
\end{lemma}

\begin{proof}
Within either prescription set, the Lipschitz inequality follows from
the bump bound. Across the two sets, the difference between the
prescribed values at $j,j'$ is
$b_x(j)+b_y(j')\le2L$. Thus the inequality holds whenever
$d_p(j,j')\ge2L$.

Suppose $d_p(j,j')<2L$ and $b_x(j)>0$. Then
$\norm{h(x)-h(j)}_\infty<K$. Since $d_p(x,y)\le1$ and $h$ is
$1$-Lipschitz, \Cref{lem:block-geometry} gives
\[
 \begin{split}
 \norm{h(y)-h(j')}_\infty
 &\le\norm{h(y)-h(x)}_\infty
       +\norm{h(x)-h(j)}_\infty
       +\norm{h(j)-h(j')}_\infty\\
 &<1+K+2L.
 \end{split}
\]
The pairs $(x,j)$ and $(y,j')$ are compatible, and their state distance
is at most $2L$. By \Cref{lem:exceptional-directions}, both
prescriptions cannot be present. The same reasoning applies when
$b_y(j')>0$. Hence whenever two states, one from each prescription set,
are less than $2L$ apart, both prescribed values are zero. This proves
the remaining Lipschitz inequalities and, by taking $j=j'$, consistency
on the intersection of the two prescription sets.

Let $P\subseteq S$ be their union and $v:P\to[-L,L]$ the resulting
$1$-Lipschitz data. If $P$ is empty, take $f_e=0$. Otherwise the
classical scalar Lipschitz extension principle \cite{McShane34}
can be realized for these finite data by the following formula,
followed by clipping:
\[
 F(z)=\min_{j\in P}\{v(j)+d_p(z,j)\},\qquad
 f_e(z)=\max\{-L,\min\{L,F(z)\}\}.
\]
For $z\in P$, the Lipschitz inequalities for $v$ give $F(z)=v(z)$.
The triangle inequality gives $F(z)\le F(w)+d_p(z,w)$ and the
same inequality with $z,w$ exchanged. Thus $F$ is $1$-Lipschitz.
Clipping to $[-L,L]$ preserves both the Lipschitz bound and the
prescribed values.
\end{proof}

\subsection{The finite anchor obstruction}

Each test compares the signed measures at the ends of its incidence.
The comparison loses variation only on exceptional states. At a good
vertex each state is exceptional in at most one of the $\Delta$
directions, so averaging these losses will give a bound independent
of the cardinality of $S$.

\begin{proof}[Proof of \Cref{thm:blocks}]
For the chosen integer $p$, \Cref{lem:block-geometry} supplies the
countably infinite block, its separation and diameter bounds, and the
finite anchor set $A_p$. Suppose that a bounded finite-rank operator
$T_0$ of norm at most $p$ has error less than $1/2$ on every anchor.
Choose $T,S$ by \Cref{lem:finite-support} and the measures $\mu_x$
by \Cref{lem:measure-bounds}. With this state set fixed, choose
$n\ge2$ and then $G=G_m$ by \Cref{lem:rare-walks}, so that the
fraction $\beta$ of bad vertices is less than $1/(20V)$. Apply
\Cref{lem:exceptional-directions} and \Cref{lem:edge-test} to these choices.

For an incidence $e:x\to y$, the test prescriptions imply
\begin{align}
 \int f_e\,d(\mu_x-\mu_y)
 &\ge\int b_x\,d\mu_x+\int b_y\,d\mu_y\notag\\
 &\quad-2L\bigl(|\mu_x|(E_e)+|\mu_y|(E_{\bar e})\bigr)\notag\\
 &\ge\frac{3L}{2}
       -2L\bigl(|\mu_x|(E_e)+|\mu_y|(E_{\bar e})\bigr).
       \label{eq:edge-lower}
\end{align}
Indeed, $f_e-b_x$ vanishes on $S\setminus E_e$ and has absolute value
at most $2L$ elsewhere; the same statement holds for $f_e+b_y$ and
$E_{\bar e}$. Integration against the variation measures bounds both
errors, regardless of the signs of the coefficients.

On the other hand, the affine pairing \eqref{eq:affine-pairing} gives
\begin{equation}\label{eq:edge-upper}
 \left|\int f_e\,d(\mu_x-\mu_y)\right|
 =\left|(T(\delta_x-\delta_y))(f_e-f_e(o))\right|
 \le2p,
\end{equation}
because $\norm{T}\le2p$, $\Lipnorm{f_e}\le1$, and
$\norm{\delta_x-\delta_y}=d_p(x,y)\le1$. For a loop, $x=y$ and
the left side is zero.

Choose a uniform vertex $x$ of $G$ and then a uniform incidence
$e\in\mathcal D(x)$. At a good vertex each fixed state occurs in
at most one exceptional set. Consequently
\[
 \frac1\Delta\sum_{e\in\mathcal D(x)}|\mu_x|(E_e)
 \le\frac1\Delta\sum_{j\in S}|\mu_x(\{j\})|
 \le\frac V\Delta.
\]
At a bad vertex the average is at most $V$. Therefore
\begin{equation}\label{eq:average-error}
 \E\bigl[|\mu_x|(E_e)\bigr]\le\frac V\Delta+V\beta<\frac1{10}.
\end{equation}
The $\Delta$-regularity of $G$ makes this sampling uniform on all oriented
incidences.
Reversal preserves that uniform distribution, exchanging the two
incidences of a loop as well. Hence
$\E\bigl[|\mu_y|(E_{\bar e})\bigr]<1/10$ too. Taking expectations in
\eqref{eq:edge-lower} and using \eqref{eq:edge-upper} forces
\[
 2p>\frac{3L}{2}-\frac{4L}{10}
    =\frac{11L}{10}=11p,
\]
a contradiction.

The construction fixed $p$, all constants, the entire graph sequence,
$M_p$, and $A_p$ before $T_0$ was considered. Only after $T_0$ produced
the finite support $S$ were $n$ and then $m$ chosen. Thus the same finite
anchor set defeats every operator in the stated class, proving
\Cref{thm:blocks}.
\end{proof}

\Cref{thm:blocks} now supplies the hypothesis used in \Cref{sec:reduction}.
The wedge construction there therefore proves \Cref{thm:main}, and the
coefficient identification proves \Cref{cor:renorming}.
\Cref{sec:nonlinear} then gives failure of nonlinear approximability for
the same space $\F(M)$.

\begin{remark}[A sufficient criterion for the anchor obstruction]\label{rem:abstract-obstruction}
The preceding proof gives the same finite-anchor obstruction for any data
satisfying the following assumptions. Fix an integer $p\ge1$
and retain the parameter choices
$L,R,K,D,B,V,\Delta,Q,H,k$ of \Cref{sec:blocks}, including
\eqref{eq:parameters}, and the incidence convention of that section.
\begin{enumerate}[label=(\roman*)]
\item $(N,d,o)$ is a countable pointed metric space of separation at
least one and diameter at most $B$, and $A\subseteq N$ is finite
and contains $o$.
\item The $G_m$ are nonempty finite $\Delta$-regular multigraphs with disjoint
vertex sets in $N$ and colors in $\{1,\ldots,k\}$. Edge endpoints
have $d$-distance at most one. For every fixed integer $r\ge0$, the
empirical distributions of their induced rooted radius-$r$ colored
balls converge to the corresponding distribution on the $\Delta$-regular
tree with independent uniform colors, as in \Cref{lem:local-limit}.
\item The map $h:N\to\{0,\ldots,D\}^k$ is $1$-Lipschitz for the sup
norm. Every graph vertex $x$ of color $i$ satisfies $h_i(x)=0$ and
has an anchor $a(x)\in A$ with $h(a(x))=h(x)$ and $d(x,a(x))\le R$.
\item Every closed $d$-ball of radius $2L$ has at most $Q$ points,
and for every $t\in N$ at most $H$ coordinates of $h(t)$ are strictly below $D$.
\end{enumerate}
If these data are fixed before $T_0$, every bounded finite-rank real
linear operator $T_0:\F(N)\to\F(N)$ with $\norm{T_0}\le p$ satisfies
\[
 \max_{a\in A}\norm{T_0\delta_a-\delta_a}\ge\frac12.
\]
Indeed, the support perturbation and signed-measure estimates use only
the stated metric, anchor, and label bounds, while the rare-walk,
continuation, scalar-test, and averaging steps use only those bounds,
the stated graph limit, and the incidence convention.
\end{remark}
